% Guia de notação, teoria e comparação de métodos do projeto vida-grupo-pricing.
% Compilar com XeLaTeX (duas vezes, para o sumário):  xelatex notacao.tex
\documentclass[11pt,a4paper]{article}

\IfFileExists{mathpazo.sty}{\usepackage[sc]{mathpazo}}{}
\usepackage{fontspec}
\IfFontExistsTF{TeX Gyre Pagella}{\setmainfont{TeX Gyre Pagella}}{}
\IfFontExistsTF{DejaVu Sans Mono}{\setmonofont[Scale=0.85]{DejaVu Sans Mono}}{}
\usepackage{polyglossia}
% português com hifenização, se os padrões estiverem instalados (TeX Live e MiKTeX completos têm)
\makeatletter
\ifcsname l@portuges\endcsname
  \setdefaultlanguage[variant=brazilian]{portuguese}
\else
  \setdefaultlanguage{english}
\fi
\makeatother

\usepackage[a4paper,margin=2.3cm]{geometry}
\usepackage{amsmath,amssymb,mathtools}
\usepackage{booktabs,tabularx,array,longtable}
\usepackage[shortlabels]{enumitem}
\usepackage{xcolor}
\usepackage{titlesec}
\usepackage{fancyhdr}
\usepackage{microtype}
\usepackage[hidelinks]{hyperref}
\hypersetup{colorlinks=true,linkcolor=azul,urlcolor=azul}

\definecolor{azul}{HTML}{1F4E79}
\titleformat{\section}{\Large\bfseries\color{azul}}{\thesection}{0.8em}{}
\titleformat{\subsection}{\large\bfseries}{\thesubsection}{0.6em}{}
\setlist{itemsep=2pt,topsep=4pt}
\newcolumntype{L}{>{\raggedright\arraybackslash}X}
\renewcommand{\arraystretch}{1.25}
\pagestyle{fancy}\fancyhf{}\renewcommand{\headrulewidth}{0pt}
\fancyfoot[C]{\small\thepage}\fancyfoot[R]{\small\color{gray}vida-grupo-pricing}

\newcommand{\onde}[1]{\par\noindent{\small\color{gray}Onde aparece: #1}\par\medskip}
\newcommand{\razaoAE}{\mathrm{A/E}}
\newcommand{\nef}{n_{\mathrm{ef}}}

\title{\color{azul}\textbf{Notação, teoria e comparação de métodos}\\[4pt]\large Precificação e renovação de Vida em Grupo}
\author{Projeto \texttt{vida-grupo-pricing}}
\date{}

\begin{document}
\maketitle
\thispagestyle{empty}

\noindent Guia de referência dos símbolos e métodos do projeto, na ordem do pipeline. Cada seção diz o que é, a
fórmula e onde aparece (notebook e módulo de \texttt{src/vgpricing}). A Seção~\ref{sec:comparacao} compara os métodos
(vantagens, desvantagens e quando usar).

\renewcommand{\contentsname}{Sumário}
\tableofcontents
\newpage

% ============================================================================================
\section{Símbolos gerais}

\begin{tabularx}{\textwidth}{@{}lL@{}}
\toprule
Símbolo & Significado \\
\midrule
$x$ & idade completa da vida \\
$q_x$ & probabilidade de morte em um ano na idade $x$ (tábua BR-EMSmt 2026, por sexo) \\
$t$ & fração do ano em que a vida esteve exposta (dias $\div$ 365,25) \\
$C$, $C_{\mathrm{ref}}$ & capital segurado; capital de referência (R\$ 50 mil), que só normaliza o fator de capital \\
$c$ & célula de modelagem (grupo $\times$ faixa etária $\times$ sexo $\times$ faixa de capital) \\
$g$, $i$, $j$ & grupo (estipulante); na credibilidade, $i$ indexa grupos e $j$ períodos \\
$N$, $E$, $\mu$ & óbitos observados; esperados pela tábua; esperados por um modelo \\
$A$ & observado (contagem ou valor), com IBNR quando indicado \\
$P$ & nível de pooling do grupo \\
$Z$, $k$ & credibilidade (peso da experiência própria) e constante de credibilidade \\
\bottomrule
\end{tabularx}

% ============================================================================================
\section{Exposição e A/E}
\onde{01 e 02 --- \texttt{exposicao}, \texttt{estudo}}

\begin{equation}
E = \sum_{\text{vidas}} q_x\, t, \qquad E^{(v)} = \sum_{\text{vidas}} q_x\, t\, C, \qquad \razaoAE = \frac{N}{E}.
\end{equation}

\textbf{Vidas-ano} são $\sum t$ e \textbf{capital-ano}, $\sum tC$.

\paragraph{Número efetivo de óbitos.} Com capitais diferentes, a variância do valor observado é $\sum C^2 q t$
(Poisson composta), e o A/E em valor tem a precisão de
\begin{equation}
\nef = \frac{\left(\sum C q t\right)^2}{\sum C^2 q t} \;\le\; E,
\end{equation}
com igualdade quando os capitais são iguais. É o peso da credibilidade em valor.

\paragraph{Intervalo exato de Poisson (Garwood).} Para a média de uma contagem $N$, ao nível $1-\alpha$:
\begin{equation}
\left[\tfrac12\,\chi^2_{\alpha/2}(2N),\ \tfrac12\,\chi^2_{1-\alpha/2}(2N+2)\right],
\end{equation}
dividido por $E$ para o A/E. Aproximação: $\razaoAE \pm 1{,}96/\sqrt{N}$. Com 30 óbitos, cerca de $\pm 36\%$: daí o
mínimo de 30 esperados para um nível ter fator próprio.

% ============================================================================================
\section{GLM de Poisson com offset}
\onde{03 e 04 --- \texttt{glmm} (R \texttt{stats::glm})}

\begin{equation}
N_c \sim \mathrm{Poisson}(\mu_c), \qquad \log \mu_c = \underbrace{\log E_c}_{\text{offset}} + \beta_0 + \sum_k \beta_k x_{ck}.
\end{equation}
O offset fixa a tábua: cada $e^{\beta_k}$ é um fator multiplicativo sobre ela.

\begin{itemize}
\item \textbf{Deviance:} $D = 2\sum_c \bigl[N_c \log(N_c/\mu_c) - (N_c - \mu_c)\bigr]$ (com $N_c = 0$, o termo vale
  $\mu_c$). Menor é melhor; é a medida de erro adequada a contagens de eventos raros.
\item \textbf{AIC} $= 2p - 2\log L$, com $p$ parâmetros; só compara modelos ajustados aos mesmos dados.
\item \textbf{ANOVA sequencial:} deviance explicada por cada variável, na ordem da fórmula (tabela da SOA, seção 4.1).
\item \textbf{Dispersão de Pearson} $= \sum (N-\mu)^2/\mu \,\big/\, (n-p)$: perto de 1, a Poisson descreve a variância.
\item \textbf{Razão de verossimilhança:} $LR = 2(\log L_1 - \log L_0) \sim \chi^2_{gl}$. Para uma variância de efeito
  aleatório (fronteira do espaço de parâmetros), $p = \tfrac12 P(\chi^2_1 > LR)$.
\item \textbf{Correção de Holm} para $m$ testes: com os p-valores ordenados,
  $p^{\mathrm{adj}}_{(k)} = \max_{j \le k} \min\{1,\ (m-j+1)\,p_{(j)}\}$.
\end{itemize}

% ============================================================================================
\section{Modelo misto (GLMM)}
\onde{04 --- \texttt{glmm} (R \texttt{glmmTMB})}

\begin{equation}
\log \mu_c = \log E_c + \mathbf{x}_c^\top\boldsymbol\beta + u_{g(c)} + v_{d(c)}, \qquad
u_g \sim N(0, \sigma^2_{\mathrm{est}}),\quad v_d \sim N(0, \sigma^2_{\mathrm{div}}).
\end{equation}

\begin{itemize}
\item O efeito do estipulante $u_g$ capta a mortalidade própria do grupo; a divisão CNAE $v_d$, dentro da seção, dá
  credibilidade a cada divisão em torno da sua seção (Ohlsson, 2008).
\item \textbf{Grupo novo} (sem $u_g$ observado): $\mathrm{E}[e^{u}] = e^{\sigma^2/2}$, não 1.
\item \textbf{Ajuste singular:} variância estimada abaixo de $10^{-8}$; o efeito é retirado e o modelo, reajustado.
\end{itemize}

% ============================================================================================
\section{Penalização: elastic net, ridge e priori}
\onde{03 e 04 --- \texttt{elasticnet} (R \texttt{glmnet})}

\begin{equation}
\hat{\boldsymbol\beta} = \arg\min_{\boldsymbol\beta}\; -\frac1n \log L(\boldsymbol\beta)
  + \lambda\left[\frac{1-\alpha}{2}\,\lVert\boldsymbol\beta\rVert_2^2 + \alpha\,\lVert\boldsymbol\beta\rVert_1\right].
\end{equation}

\begin{itemize}
\item $\alpha = 1$: LASSO (zera coeficientes); $\alpha = 0$: ridge (encolhe todos). $\lambda$ por validação cruzada
  (\texttt{lambda.min}, o de menor erro).
\item \textbf{Ridge é uma priori normal.} Maximizar a verossimilhança penalizada por $\lambda\lVert\beta\rVert^2$ é
  achar a moda da posteriori com $\beta \sim N(0, \tau^2)$, $\lambda \propto 1/\tau^2$. LASSO corresponde a uma
  priori de Laplace.
\item \textbf{Priori de um estudo de mercado.} Com a forma do estudo no offset,
  \begin{equation}
  \log\mu = \log E + \beta^{\mathrm{SOA}}\log\tfrac{C}{C_{\mathrm{ref}}} + \log f^{\mathrm{SOA}}_s
          + b_0 + \delta\,\log\tfrac{C}{C_{\mathrm{ref}}} + \gamma_s ,
  \end{equation}
  a penalização encolhe $(\delta, \gamma_s)$ para zero, isto é, para o estudo: $\lambda$ alto mantém o estudo; dados
  abundantes seguem a carteira. A SOA (2024, seção 4.3.1) liga $\lambda$ ao $k$ de Bühlmann.
\end{itemize}

\subsection{AGLM}
\onde{03 e 04 --- \texttt{aglm} sobre \texttt{glmnet}}

GLM regularizado com codificações que captam não linearidade sem perder a forma de tabela (Fujita, Tanaka, Kondo e
Iwasawa, 2020).
\begin{itemize}
\item \textbf{Variáveis O} (dummies ordinais): para níveis ordenados $1,\dots,m$, $d^O_j(x) = \mathbb{1}\{x < j\}$,
  $j = 2,\dots,m$. O coeficiente $\beta_j$ é o salto entre os níveis $j-1$ e $j$; com LASSO, $\beta_j = 0$ funde as
  duas faixas. Equivale ao \emph{fused lasso}, $\lambda\sum_j |\gamma_j - \gamma_{j-1}|$ sobre os efeitos de nível.
\item \textbf{Variáveis L}: com nós $b_1 < \dots < b_{m-1}$, $l_j(x) = |x - b_j|$, mais o próprio $x$. A contribuição
  $\sum_j \beta_j l_j(x)$ é linear por partes e contínua; $\beta_j$ é a mudança de inclinação no nó $b_j$.
\item No projeto: capital em variáveis L (9 nós nos quantis do log), faixa etária e porte em O, seção e sexo em dummies
  comuns, LASSO e a forma da SOA no offset.
\end{itemize}

\subsection{GAM}
\onde{03 e 04 --- \texttt{gam} (R \texttt{mgcv})}

\begin{equation}
\log\mu = \log E + \beta_0 + f(\log C) + \mathbf{x}^\top\boldsymbol\beta, \qquad f(z) = \sum_k \theta_k B_k(z),
\end{equation}
com $B_k$ uma base de splines e penalização de rugosidade $\lambda\int f''(z)^2\,dz$, escolhida por REML. Os
\textbf{graus de liberdade efetivos} (edf) medem a complexidade usada: edf $\approx 1$ é uma reta (no log, uma
elasticidade constante). A penalização é uma priori gaussiana sobre a rugosidade (Wood, 2017).

\subsection{LightGBM, SHAP e interações}
\onde{03 e 04 --- \texttt{diagnostico}}

\begin{itemize}
\item \textbf{Gradient boosting} com perda Poisson e valor inicial $\log E$: soma árvores que corrigem o log da taxa.
\item \textbf{SHAP:} decomposição aditiva $\log\mu(\mathbf{x}) = \log\mu_0 + \sum_j \phi_j(\mathbf{x})$; importância
  $=$ média de $|\phi_j|$.
\item \textbf{Tabela de interações (estilo EIX):} cada divisão de uma árvore cujo nó pai divide por outra variável forma
  um par, que acumula o ganho do filho (\emph{sumGain}), as observações do filho (\emph{sumCover}) e a frequência.
\end{itemize}

% ============================================================================================
\section{Credibilidade}
\onde{05 --- \texttt{renovacao} (R \texttt{actuar})}

\paragraph{Flutuação limitada (prática de mercado).} $Z = \min\bigl(1, \sqrt{n/n_{\mathrm{full}}}\bigr)$, com $n$ em
vidas-ano, $n_{\mathrm{full}} = 20\,000$ e $Z = 0$ abaixo de 500.

\paragraph{Bühlmann-Straub.} Grupo $i$, período $j$, razão observada $X_{ij}$ (A/E em valor contra a taxa manual) e
peso $w_{ij}$ ($\nef$). Com $w_i = \sum_j w_{ij}$, $\bar X_i = \sum_j w_{ij} X_{ij}/w_i$ e
$\bar X = \sum_i w_i \bar X_i / w$:
\begin{align}
\hat s^2 &= \frac{\sum_i \sum_j w_{ij}\,(X_{ij} - \bar X_i)^2}{\sum_i (n_i - 1)}, &
\hat a &= \frac{\sum_i w_i\,(\bar X_i - \bar X)^2 - (I-1)\,\hat s^2}{w - \sum_i w_i^2/w}, \\[4pt]
k &= \frac{\hat s^2}{\hat a}, &
Z_i &= \frac{w_i}{w_i + k}, \qquad \hat X_i = Z_i\,\bar X_i + (1 - Z_i)\,\mu .
\end{align}
$s^2$ é a variância dentro do grupo (ruído) e $a$ a variância entre grupos (sinal). Se $\hat a \le 0$, ou se não há
dois grupos com dois períodos, $Z = 0$. Na renovação, o complemento $\mu$ é a própria taxa manual (A/E $= 1$). É a
melhor previsão linear (Bayes linear) do risco do grupo.

% ============================================================================================
\section{Pooling e IBNR}
\onde{05 --- \texttt{renovacao}, \texttt{ibnr}}

\begin{itemize}
\item \textbf{Nível de pooling:} $P = \min\{\text{retenção},\ \max(\text{piso},\ m\,\bar C)\}$, com $\bar C$ o capital
  médio do grupo. Na experiência entra $\min(C, P)$; o excesso esperado vira encargo cobrado de todos:
  $\sum q f \max(C - P, 0)$.
\item \textbf{Chain ladder mensal} (carteira): fatores de desenvolvimento $f_j = \sum_i S_{i,j+1} / \sum_i S_{i,j}$,
  $\mathrm{CDF}_j = \prod_{k \ge j} f_k$ e proporção avisada $p_j = 1/\mathrm{CDF}_j$.
\item \textbf{Bornhuetter-Ferguson por grupo:} $\mathrm{IBNR} = \sum_{\text{meses}} E^{\text{apriori}}_m\,(1 - p_{a(m)})$,
  com $a(m)$ a idade do mês de ocorrência na avaliação; estável para grupos com poucos sinistros.
\item \textbf{Fator de desenvolvimento da taxa manual:} média de $1/p_a$ sobre os sinistros avisados da experiência
  (cada sinistro avisado com idade $a$ representa $1/p_a$ sinistros finais).
\end{itemize}

% ============================================================================================
\section{Reajuste e cotação}
\onde{05 e 06 --- \texttt{renovacao}, \texttt{cotador}}

\begin{gather}
\text{prêmio puro} = \sum q\, f_{\mathrm{manual}} \min(C,P)\cdot \text{fator} + \text{encargo}, \qquad
\text{taxa comercial} = \text{taxa pura}\cdot\frac{1 + \text{margem}}{1 - \sum\text{carregamentos}}, \\[4pt]
\text{reajuste} = \operatorname{clip}\!\left(\frac{\text{taxa comercial}}{\text{taxa vigente}} - 1,\ \text{piso},\ \text{teto}\right),
\end{gather}
com $f_{\mathrm{manual}} =$ A/E base $\times$ fator de seção $\times$ fator de divisão $\times (C/C_{\mathrm{ref}})^{\beta_{\mathrm{cap}}}$.
Na renovação, $\text{fator} = Z\,\razaoAE + (1-Z)$; na cotação de grupo novo sem histórico, $\text{fator} = 1$.

\paragraph{Cotação com o histórico do grupo.} A subscrição recebe, por ano $j$ do histórico, os sinistros avisados até
a data da cotação, $S_j$ (limitados ao pooling), e o capital $K_j$. Com $p_j$ a proporção do ano $j$ que já deveria
estar avisada (padrão de aviso da carteira) e $c$ o custo esperado anual do censo na taxa manual,
\begin{equation}
\razaoAE_{\mathrm{hist}} = \frac{\sum_j S_j/p_j}{\sum_j c\,K_j/K}, \qquad
\nef^{\mathrm{hist}} = \nef \sum_j \frac{K_j}{K}, \qquad
Z = \frac{\nef^{\mathrm{hist}}}{\nef^{\mathrm{hist}} + k},
\end{equation}
com o $k$ de Bühlmann-Straub da carteira. Sem dividir por $p_j$, o último ano (incompleto) puxa a taxa para baixo.

% ============================================================================================
\section{Validação}
\onde{03, 04, 05 e 06 --- \texttt{avaliacao}, \texttt{backtest}}

\begin{itemize}
\item \textbf{Teste fora do tempo:} ajusta-se com dados até uma data e avalia-se no período seguinte; o teste só é
  usado no fim (CAS, seção 4.3).
\item \textbf{Bootstrap pareado:} reamostram-se as unidades (estipulantes, seguradoras) e soma-se a diferença de
  deviance entre dois modelos; o IC de 95\% são os percentis 2,5 e 97,5. ``Significativa'' $=$ IC sem o zero.
\item \textbf{Gráfico de quantis:} ordena-se pela relatividade prevista $\mu/E$, em faixas de igual $E$ (CAS 7.2.1).
\item \textbf{Double lift:} ordena-se pela razão entre dois modelos; vence o que acompanha o observado (CAS 7.2.2).
\item \textbf{Lorenz e Gini:} do menor para o maior risco previsto, $x$ é a parcela acumulada de $E$ e $y$ a de $N$;
  Gini $= 2\times$ a área entre a diagonal e a curva (CAS 7.2.4).
\item \textbf{A/M:} observado $\div$ modelo por nível; desvios além de $\pm 5\%$ a $\pm 10\%$ indicam ajuste ruim
  (SOA 2.4.5); aqui também com IC exato.
\item \textbf{Gráfico de sinistralidade:} ordena-se pelo custo previsto $\div$ prêmio atual (CAS 7.2.3).
\item \textbf{Estabilidade:} bootstrap por estipulante dos fatores da tarifa (CAS 6.4).
\item \textbf{Cotação com histórico:} grupos novos simulados como chegam à subscrição (censo e histórico avisado),
  comparados com o ano seguinte, que a cotação não vê.
\item \textbf{Regra de decisão:} menor deviance no teste; empatados (IC contém zero) resolvidos por uma ordem de
  preferência fixada antes.
\end{itemize}

% ============================================================================================
\section{Monitoramento}
\onde{06 --- \texttt{monitoramento}}

\begin{itemize}
\item \textbf{Limites da carteira:} sob ``a taxa está certa'', $N \sim \mathrm{Poisson}(E)$; os limites do A/E são os
  quantis de $\mathrm{Poisson}(E)$ divididos por $E$, a 95\% e 99,8\% (funnel plot, Spiegelhalter, 2005).
\item \textbf{Limites dos grupos:} com mortalidade própria $u \sim \mathrm{Gama}$ de média 1 e variância $\tau^2$,
  $N \sim$ Binomial Negativa com média $E$ e variância $E + \tau^2 E^2$. Estimador de momentos:
  $\hat\tau^2 = \max\bigl\{0,\ \sum[(N-E)^2 - N] \big/ \sum E^2\bigr\}$.
\item \textbf{IBNR nos meses recentes:} $N^{\mathrm{aj}} = N/p_a$; os meses com $a \le 3$ ficam fora do A/E móvel.
\end{itemize}

% ============================================================================================
\section{Dados faltantes}
\onde{01 --- \texttt{exposicao}, \texttt{validacao}, \texttt{manual}}

\begin{itemize}
\item \textbf{Sexo ausente:} com $p$ a proporção de homens do perfil,
  \begin{equation}
  \mathrm{E}[q] = p\,q^M_x + (1-p)\,q^F_x,
  \end{equation}
  exata para o número esperado de óbitos, que é linear em $p$. $p$ vem do nível mais fino com dados: grupo e faixa de
  10 anos, grupo, seção CNAE e faixa, carteira; para cotar um grupo sem sexo, a tabela de fatores guarda a proporção
  da carteira por seção e faixa etária. A média 50/50 só é certa se $p = 0{,}5$.
\item \textbf{Ausência informativa:} se a falta de sexo se concentra em certos grupos ou atividades, a ausência carrega
  informação; por isso ``ND'' é um nível próprio nas células, testado como fator.
\item \textbf{Nascimento ausente:} imputação pela idade de entrada mediana do grupo, com indicador. Excluir a vida
  retira exposição e enviesa o A/E.
\item \textbf{Gradient boosting} aceita ausentes nativamente, mas de forma implícita: o modelo pode usar ``ausente''
  como substituto de outra coisa. Além disso, o sexo entra pela tábua (no offset). A ASOP 23 pede tratamento explícito
  e documentado.
\end{itemize}

% ============================================================================================
\section{Comparação dos métodos}\label{sec:comparacao}

{\small
\begin{longtable}{@{}>{\raggedright\arraybackslash}p{0.17\textwidth}>{\raggedright\arraybackslash}p{0.27\textwidth}%
  >{\raggedright\arraybackslash}p{0.27\textwidth}>{\raggedright\arraybackslash}p{0.22\textwidth}@{}}
\toprule
Método & Vantagens & Desvantagens & Quando usar \\
\midrule\endhead
\textbf{GLM de Poisson com offset} & interpretável fator a fator; erro-padrão; padrão regulatório e de mercado; rápido
  & forma fixada pelo modelador; não encolhe níveis raros & tarifa, com volume para os fatores \\
\textbf{GLMM} & credibilidade automática para níveis e grupos; separa ruído de heterogeneidade
  & efeitos fixos condicionais ao grupo (como taxa de grupo novo, podem prever pior); ajuste singular
  & estimar a variância entre grupos; atividades com muitos níveis \\
\textbf{Ridge / elastic net} & encolhe níveis raros; aceita priori via offset; LASSO seleciona variáveis
  & sem erro-padrão (estabilidade por bootstrap); $\lambda$ depende da validação cruzada
  & pouca experiência; muitos níveis; usar estudo de mercado \\
\textbf{AGLM} & não linearidade e fusão de faixas sem escolher bins à mão; continua sendo tabela
  & curvas podem não ser monótonas; mais parâmetros para explicar & contínuas ou ordenadas com forma desconhecida \\
\textbf{GAM} & forma suave e estimada; edf mostra se a forma linear basta & curva precisa ser tabulada para tarifa
  & diagnóstico da forma de variáveis contínuas \\
\textbf{LightGBM} & capta interações e não linearidades sem especificar; aceita ausentes
  & não explicável fator a fator; exige regularização forte com eventos raros
  & referência de desempenho; busca de estrutura faltante \\
\textbf{Credibilidade de mercado} & simples, conhecida no mercado & limiares arbitrários; ignora concentração de capital
  & comparação (\emph{champion}) \\
\textbf{Bühlmann-Straub} & $k$ estimado nos dados; peso pelo número efetivo; melhor previsão linear
  & supõe estrutura estável entre períodos; $k$ incerto com poucos grupos & renovação e cotação com histórico \\
\textbf{Chain ladder + BF} & padrão de mercado; BF estável para grupos pequenos & supõe padrão de aviso estável
  & IBNR da carteira, por grupo e do histórico na cotação \\
\bottomrule
\end{longtable}
}

\section*{Referências}
\begin{itemize}[leftmargin=*]
\item Bühlmann, H.; Gisler, A. (2005). \emph{A Course in Credibility Theory and its Applications}. Springer.
\item Fujita, S.; Tanaka, T.; Kondo, K.; Iwasawa, H. (2020). \emph{AGLM: A Hybrid Modeling Method of GLM and Data
  Science Techniques}.
\item Goldburd, M.; Khare, A.; Tevet, D.; Guller, D. (2025). \emph{Generalized Linear Models for Insurance Rating},
  2.\,ed. CAS Monograph 5.
\item Ohlsson, E. (2008). Combining generalized linear models and credibility models in practice.
  \emph{Scandinavian Actuarial Journal}.
\item Society of Actuaries (2024). \emph{A Predictive Analytics Framework: A Reusable Workflow for Experience Analysis}.
\item Spiegelhalter, D. (2005). Funnel plots for comparing institutional performance. \emph{Statistics in Medicine}.
\item Wood, S. (2017). \emph{Generalized Additive Models: An Introduction with R}, 2.\,ed. CRC Press.
\end{itemize}

\end{document}
